\section{The linear-readout floor}
\label{sec:floor}
% ============================================================================

\begin{theorem}[Unit-diagonal cross-Gramian Welch floor]
\label{thm:floor}
Let $F\ge2$, let $\Phi\in\R^{d\times F}$ and $G\in\R^{F\times d}$, and set $M=G\Phi$. Suppose
$M_{ii}=1$ for all $i\in[F]$. Then
\begin{equation}
  \sum_{i\neq j}|M_{ij}|^2\;\ge\;\frac{F(F-d)}{d}.
  \label{eq:floorsum}
\end{equation}
Consequently, if $F>d$,
\begin{equation}
  \frac{1}{F(F-1)}\sum_{i\neq j}|M_{ij}|^2\;\ge\;W(F,d):=\frac{F-d}{d(F-1)},
  \label{eq:floor}
\end{equation}
and hence $\max_{i\neq j}|M_{ij}|\ge\sqrt{W(F,d)}$. In particular, if $F/d\to\infty$ then
$\max_{i\neq j}|M_{ij}|\ge(1-o(1))\,d^{-1/2}$.
\end{theorem}

\begin{proof}
Since $M=G\Phi$ with $G\in\R^{F\times d}$ and $\Phi\in\R^{d\times F}$ we have
$\operatorname{rank}(M)\le d$, and the unit diagonal gives $\operatorname{tr}(M)=F$. For any
matrix of rank at most $d$, $|\operatorname{tr}(M)|\le\nucnorm{M}\le\sqrt d\,\norm{M}_F$
(Appendix~\ref{app:nuclear}), so $F^2\le d\,\norm{M}_F^2$. Expanding
$\norm{M}_F^2=F+\sum_{i\neq j}|M_{ij}|^2$ gives~\eqref{eq:floorsum}, and dividing by $F(F-1)$
gives~\eqref{eq:floor}. For the maximum, note that
$\max_{i\neq j}|M_{ij}|^2\ge\frac{1}{F(F-1)}\sum_{i\neq j}|M_{ij}|^2$, and take square roots.
\end{proof}

The argument is the standard rank--trace mechanism, cf.~\cite{datta2012geometry}; we give it
in full only for self-containedness. The next corollary is what the diagnostic actually uses,
because it removes the normalisation hypothesis.

\begin{corollary}[Gain-normalised form]
\label{cor:calibfree}
Let $F>d$ and let $M=G\Phi$ satisfy $\operatorname{rank}(M)\le d$ and $M_{ii}\neq0$ for all
$i$. Then
\[
  \max_{i\neq j}\frac{|M_{ij}|}{|M_{ii}|}\;\ge\;\sqrt{\frac{F-d}{d(F-1)}}.
\]
\end{corollary}

\begin{proof}
Apply Theorem~\ref{thm:floor} to $D^{-1}M$ with $D=\operatorname{diag}(M_{11},\dots,M_{FF})$:
this matrix has rank at most $d$, unit diagonal, and entries $M_{ij}/M_{ii}$.
\end{proof}

\begin{remark}[Why the gain normalisation is essential]
\label{rem:calib}
Without a diagonal normalisation the absolute off-diagonal entries of $G\Phi$ can be made
arbitrarily small by shrinking $G$, and the floor appears to fail. A concrete instance: take
$\Phi$ to be the three unit vectors at $120^{\circ}$ in $\R^2$ --- an equiangular tight frame
--- and $G=D\Phi^{\top}$ with $D=\operatorname{diag}(1/10,1/20,1/20)$. Then $G\Phi$ has
diagonal $(1/10,1/20,1/20)$ and maximum off-diagonal entry $1/20$, ten times \emph{below}
$\sqrt{(F-d)/(d(F-1))}=1/2$. There is no contradiction: this interface reads out its own
coordinates with gains far below one, and after recalibrating each row to unit gain the
maximum off-diagonal entry becomes exactly $1/2$ --- the floor, attained with equality.
Corollary~\ref{cor:calibfree} isolates the scale-invariant content: what is bounded below is
the interference-to-gain ratio. This example is verified in the test suite, and
Section~\ref{sec:res-e4} shows that a gradient optimiser left free of the constraint exploits
exactly this degree of freedom.
\end{remark}

\begin{proposition}[Equality within the tied family]
\label{prop:tight}
Let $\Phi\in\R^{d\times F}$ have unit-norm columns and let $G=\Phi^{\top}$, so that
$M=\Phi^{\top}\Phi$ has unit diagonal. Then $\sum_{i\neq j}|M_{ij}|^2\ge F(F-d)/d$, with
equality \emph{if and only if} $\Phi\Phi^{\top}=(F/d)I_d$, that is, iff $\Phi$ is a unit-norm
tight frame.
\end{proposition}

\begin{proof}
The diagonal entries are $\norm{\phi_i}^2=1$. By cyclicity,
$\norm{M}_F^2=\operatorname{tr}((\Phi\Phi^{\top})^2)=\norm{\Sigma}_F^2$ with
$\Sigma=\Phi\Phi^{\top}\succeq0$ and $\operatorname{tr}\Sigma=F$. Cauchy--Schwarz on the
eigenvalues of $\Sigma$ gives $\norm{\Sigma}_F^2\ge(\operatorname{tr}\Sigma)^2/d=F^2/d$, with
equality iff all $d$ eigenvalues coincide, i.e.\ $\Sigma=(F/d)I_d$. Subtracting the diagonal
contribution $F$ gives the claim.
\end{proof}

Within this family --- unit-norm columns read out by their own transpose --- equality therefore
characterises tight frames, consistent
with~\cite{waldron2003generalized,datta2012geometry}. We stress the scope: it is a statement
about the tied Gram matrix, not about every rank-$d$ unit-diagonal interface. Other readouts
attain the same floor under different conditions, and the next proposition identifies the
condition for the pseudoinverse. Whether learned codes find such configurations is an empirical
question, answered in Section~\ref{sec:res-e4}.

\begin{proposition}[The pseudoinverse ratio is a leverage statistic]
\label{prop:leverage}
Let $\Phi\in\R^{d\times F}$ have full row rank, let $G=\Phi^{+}$, and write
$P=\Phi^{+}\Phi$ for the orthogonal projector onto the row space, with leverage scores
$h_i=P_{ii}\in(0,1]$. Then the gain-normalised interface $\widetilde M=D^{-1}P$ satisfies
\[
  \sum_{i\neq j}\widetilde M_{ij}^{2}\;=\;\sum_{i=1}^{F}\frac1{h_i}-F ,
  \qquad\text{so}\qquad
  r=\frac{d}{F(F-d)}\left(\sum_{i=1}^{F}\frac1{h_i}-F\right).
\]
Since $\sum_i h_i=\operatorname{tr}P=d$, Cauchy--Schwarz gives $\sum_i h_i^{-1}\ge F^2/d$ with
equality iff $h_i=d/F$ for every $i$. Hence $r\ge1$, and $r=1$ exactly when leverage is
equalised across features.
\end{proposition}

\begin{proof}
$P$ is symmetric idempotent, so row $i$ obeys $\sum_j P_{ij}^2=P_{ii}=h_i$ and $D_{ii}=h_i$.
Dividing row $i$ by $h_i$ and removing the unit diagonal entry leaves
$(h_i-h_i^2)/h_i^2=h_i^{-1}-1$; summing over $i$ gives the identity. The inequality is
Cauchy--Schwarz applied to $\sum_i h_i\cdot\sum_i h_i^{-1}\ge F^2$, and dividing by $F(F-1)$
and by $W(F,d)=(F-d)/(d(F-1))$ turns it into $r\ge1$.
\end{proof}

This is worth stating because it deflates a natural over-reading of a ratio close to one under
the pseudoinverse readout. Attaining the floor with $G=\Phi^{+}$ says that the code spreads its
features evenly over the row space --- nothing more. It is a statement about the geometry of the
code, not about whether the code is a good representation for any task; a code can have
perfectly equalised leverage and still be useless. We use this to scope the empirical claims of
Section~\ref{sec:res-e5}.

\subsection{The floor of the code itself}
\label{sec:codefloor}

Theorem~\ref{thm:floor} bounds every rank-$d$ unit-diagonal interface, so it is a statement
about the \emph{ensemble} of codes: no particular $\Phi$ need be able to attain it. That makes a
measured ratio hard to read. If a trained code sits $0.06\%$ above $W(F,d)$, is the readout
excellent, or is the floor simply unreachable for this code and the readout merely adequate? The
question is answered by computing the floor of the code itself, which is available in closed
form.

\begin{theorem}[Code-specific analog optimum]
\label{thm:codefloor}
Let $\Phi\in\R^{d\times F}$ have full row rank, let $\Sigma=\Phi\Phi^{\top}\succ0$ be its frame
operator, and let $h_i=\phi_i^{\top}\Sigma^{-1}\phi_i$ be the leverage of feature $i$. Then for
every $i$
\begin{equation}
  \min_{g\in\R^d}\Big\{\sum_{j\neq i}(g^{\top}\phi_j)^2 \;:\; g^{\top}\phi_i=1\Big\}
  \;=\;\frac{1}{h_i}-1,
  \label{eq:capon}
\end{equation}
attained uniquely at $g_i^{\star}=\Sigma^{-1}\phi_i/h_i$. Stacking the $F$ minimisers as rows gives
$G^{\star}=D_h^{-1}\Phi^{+}$ with $D_h=\operatorname{diag}(h_1,\dots,h_F)$: the row-wise optimum is
the Moore--Penrose pseudoinverse \emph{after} gain calibration, and coincides with $\Phi^{+}$ itself
only when every $h_i=1$. Consequently the smallest mean-square off-diagonal cross-talk this code
admits is
\begin{equation}
  W_{\star}(\Phi)\;:=\;\frac{1}{F(F-1)}\sum_{i=1}^{F}\Big(\frac1{h_i}-1\Big)
  \;=\;\frac{\sum_i h_i^{-1}-F}{F(F-1)}\;\ge\;W(F,d),
  \label{eq:codefloor}
\end{equation}
with equality if and only if $h_i=d/F$ for every $i$.
\end{theorem}

\begin{proof}
Since $\sum_{j}(g^{\top}\phi_j)^2=g^{\top}\Sigma g$, the constraint $g^{\top}\phi_i=1$ turns the
objective into $g^{\top}\Sigma g-1$. Minimising a positive-definite quadratic under one linear
equality is classical: the Lagrange conditions give $g=\lambda\Sigma^{-1}\phi_i$, the constraint
fixes $\lambda=1/h_i$, and the value is $\phi_i^{\top}\Sigma^{-1}\phi_i/h_i^2=1/h_i$. Subtracting
the constrained direction's contribution, which is exactly $1$, gives~\eqref{eq:capon}; strict
convexity gives uniqueness. Summing over $i$ and dividing by $F(F-1)$
gives~\eqref{eq:codefloor}. Finally $\sum_ih_i=\operatorname{tr}(\Sigma^{-1}\Phi\Phi^{\top})=d$,
so Cauchy--Schwarz gives $\sum_ih_i^{-1}\ge F^2/d$ with equality iff the $h_i$ are all equal,
and $(F^2/d-F)/(F(F-1))=W(F,d)$.
\end{proof}

\begin{remark}[Provenance]
\label{rem:capon}
Equation~\eqref{eq:capon} is the minimum-variance distortionless-response (Capon) beamformer with
the frame operator in place of a covariance~\cite{capon1969}, and $g_i^{\star}$ is the canonical
dual frame vector rescaled to unit gain; both are classical. We claim no
novelty in the optimisation. What the diagnostic uses is the reading of $h_i$ as a leverage
score, the resulting per-code floor~\eqref{eq:codefloor}, and the exact attribution of the gap
in Proposition~\ref{prop:excess}.
\end{remark}

\begin{proposition}[The excess is leverage dispersion, exactly]
\label{prop:excess}
With $c=d/F$ the mean leverage,
\[
  \sum_{i=1}^{F}\frac1{h_i}-\frac{F^2}{d}\;=\;\sum_{i=1}^{F}\frac{(h_i-c)^2}{c^{2}h_i}\;\ge\;0 .
\]
Hence $W_{\star}(\Phi)$ exceeds $W(F,d)$ by a weighted dispersion of the leverage scores, which
vanishes precisely when leverage is uniform.
\end{proposition}

\begin{proof}
Expand the right-hand side: $c^{-2}\sum_i(h_i^2-2ch_i+c^2)/h_i
=c^{-2}\sum_ih_i-2c^{-1}F+\sum_ih_i^{-1}$. Substituting $\sum_ih_i=d$ and $c=d/F$ turns the
first two terms into $F^2/d-2F^2/d=-F^2/d$.
\end{proof}

This separates two questions that a single ratio against $W(F,d)$ conflates. Writing
$W(G,\Phi)$ for the mean-square off-diagonal cross-talk actually achieved by a readout $G$,
\begin{equation}
  \underbrace{\frac{W(G,\Phi)}{W(F,d)}}_{\text{what is usually reported}}
  \;=\;\underbrace{\frac{W_{\star}(\Phi)}{W(F,d)}}_{R_{\mathrm{geom}}:\ \text{the code}}
  \times
  \underbrace{\frac{W(G,\Phi)}{W_{\star}(\Phi)}}_{R_{\mathrm{readout}}:\ \text{the decoder}} ,
  \label{eq:decomposition}
\end{equation}
and the two factors answer different questions. $R_{\mathrm{geom}}$ asks whether the code's
leverage is spread evenly; $R_{\mathrm{readout}}\ge1$ asks whether the decoder is the best one
\emph{for that code}. In particular Proposition~\ref{prop:leverage} says that the calibrated
pseudoinverse has $R_{\mathrm{readout}}=1$ identically, so a published attainment ratio computed
with $G=\Phi^{+}$ measures $R_{\mathrm{geom}}$ and nothing else. Section~\ref{sec:res-e5} reports
both factors.

Finally, the analog level reads a state distribution through one object only, which is what lets
Section~\ref{sec:hierarchy} compare it with the affine level under a single state model.

\begin{proposition}[Distribution-weighted reconstruction error]
\label{prop:secondmoment}
Let $M=G\Phi$ have unit diagonal, let $A=M-I_F$, and let the state $b\in\R^F$ be drawn from any
law with finite second moment $C=\E[bb^{\top}]$. Then
$\E\norm{Ab}_2^2=\operatorname{tr}(A^{\top}A\,C)$. The analog level therefore depends on the
state law only through $C$: two laws with the same second moment impose the same reconstruction
error on every readout.
\end{proposition}

\begin{proof}
$\E\norm{Ab}_2^2=\E\operatorname{tr}(A^{\top}Abb^{\top})=\operatorname{tr}(A^{\top}A\,C)$ by
linearity.
\end{proof}

\begin{corollary}[No better-than-floor $\varepsilon$-linear readout]
\label{cor:impossible}
Let $F>d$ and suppose $G$ satisfies $\norm{G y_i-e_i}_\infty\le\delta$ for singleton states
$y_i$, with $\delta<1/2$. Then $\delta/(1-\delta)\ge\sqrt{(F-d)/(d(F-1))}$; in particular
$\delta=\Omega(d^{-1/2})$ when $F/d\to\infty$.
\end{corollary}

\begin{proof}
Let $M=G[y_1,\dots,y_F]$. The hypothesis gives $|M_{ii}-1|\le\delta$ and $|M_{ji}|\le\delta$
for $j\neq i$. Calibrating, $|\widetilde M_{ij}|\le\delta/(1-\delta)$, and
Theorem~\ref{thm:floor} applies to $\widetilde M$.
\end{proof}

% ============================================================================
